\documentclass[11pt]{amsart}
\usepackage{amsmath,amssymb,amsthm}

\newtheorem{theorem}{Theorem}[section]
\newtheorem{lemma}[theorem]{Lemma}
\newtheorem{proposition}[theorem]{Proposition}
\newtheorem{corollary}[theorem]{Corollary}
\newtheorem{conjecture}[theorem]{Conjecture}
\theoremstyle{definition}
\newtheorem{definition}[theorem]{Definition}
\theoremstyle{remark}
\newtheorem*{remark}{Remark}

\title{Greatest common divisors of the numbers $2^n-1$}
\author{J. P. Specimen}
\address{Centre for Conjectural Arithmetic (fictional)}
\email{j.p.specimen@example.org}
\date{}

\begin{document}

\begin{abstract}
We give an elementary proof that $\gcd(2^m-1,2^n-1)=2^{\gcd(m,n)}-1$ for all
natural numbers $m,n$, by running the Euclidean algorithm on the exponents.
We derive the usual consequences for coprimality and for prime values, and
discuss the prime divisors of $2^p-1$.
\end{abstract}

\maketitle

\section{Introduction}

The sequence $2^n-1$ is a \emph{strong divisibility sequence}: the gcd of two
of its terms is again a term, with index the gcd of the indices. This
property was studied systematically by Lucas~\cite{Luc}. We present a short
self-contained proof (Theorem~\ref{thm:main}) based on a single addition
formula (Lemma~\ref{lem:add}), and collect some consequences.

\section{The gcd identity}

\begin{definition}\label{def:a}
For $n\in\mathbb{N}$ put $a_n=2^n-1$. Thus $a_0=0$, $a_1=1$, $a_2=3$,
$a_3=7$, $a_4=15$. We use the convention $\gcd(m,0)=m$ for $m\in\mathbb{N}$.
\end{definition}

\begin{lemma}\label{lem:add}
For all $m,n\in\mathbb{N}$ we have $a_{m+n}=2^n a_m+a_n$.
\end{lemma}

\begin{proof}
$2^n(2^m-1)+(2^n-1)=2^{m+n}-1$.
\end{proof}

\begin{lemma}\label{lem:odd}
For every $n\in\mathbb{N}$ the number $a_n$ is odd.
\end{lemma}

\begin{proof}
$2^n$ is even, so $2^n-1$ is odd.
\end{proof}

\begin{lemma}\label{lem:div}
Let $m,n$ be integers with $1\le m\le n$. If $m\mid n$, then $a_m\mid a_n$.
\end{lemma}

\begin{proof}
Write $n=mq$. Then
$a_n=(2^m)^q-1=(2^m-1)\bigl(1+2^m+2^{2m}+\dots+2^{(q-1)m}\bigr)$.
\end{proof}

\begin{lemma}\label{lem:step}
Let $m\ge n\ge 0$. Then $\gcd(a_m,a_n)=\gcd(a_{m-n},a_n)$.
\end{lemma}

\begin{proof}
By Lemma~\ref{lem:add} applied to the pair $(m-n,n)$ we have
$a_m=2^m a_{m-n}+a_n$, hence $\gcd(a_m,a_n)=\gcd(2^m a_{m-n},a_n)$. Since
$a_n$ is odd by Lemma~\ref{lem:odd}, it is coprime to $2^m$, and therefore
$\gcd(2^m a_{m-n},a_n)=\gcd(a_{m-n},a_n)$.
\end{proof}

\begin{theorem}[Main theorem]\label{thm:main}
For all $m,n\in\mathbb{N}$,
\[
  \gcd(2^m-1,\,2^n-1)=2^{\gcd(m,n)}-1 .
\]
\end{theorem}

\begin{proof}
We argue by strong induction on $m+n$. If $n=0$, then
$\gcd(a_m,a_0)=\gcd(a_m,0)=a_m=a_{\gcd(m,0)}$, and the case $m=0$ is symmetric.
Otherwise we may assume $m\ge n\ge 1$ by symmetry. Then by
Lemma~\ref{lem:step} and the induction hypothesis applied to $(m-n,n)$,
\[
  \gcd(a_m,a_n)=\gcd(a_{m-n},a_n)=a_{\gcd(m-n,n)}=a_{\gcd(m,n)} .\qedhere
\]
\end{proof}

\begin{corollary}\label{cor:coprime}
For $m,n\in\mathbb{N}$, the numbers $2^m-1$ and $2^n-1$ are coprime if and
only if $\gcd(m,n)=1$.
\end{corollary}

\begin{proof}
By Theorem~\ref{thm:main} the gcd is $a_g$ with $g=\gcd(m,n)$, and $a_g=1$ if
and only if $g=1$.
\end{proof}

\begin{proposition}\label{prop:vac}
Let $m,n\ge 1$ be integers such that $m\nmid n$ but $a_m\mid a_n$. Then
$n>m$.
\end{proposition}

\begin{proof}
Since $a_n\ge 1$ and $a_m\mid a_n$ we get $a_m\le a_n$, so $m\le n$; and
$m\ne n$ because $m\nmid n$.
\end{proof}

\section{Prime values and prime divisors}

\begin{proposition}\label{prop:prime}
If $a_n$ is prime, then $n$ is prime.
\end{proposition}

\begin{proof}
As $a_0=0$ and $a_1=1$, we have $n\ge 2$. If $n=rs$ with $1<r<n$, then
$a_r\mid a_n$ by Lemma~\ref{lem:div} and $1<a_r<a_n$, so $a_n$ is composite.
\end{proof}

\begin{proposition}\label{prop:divisors}
Let $p$ be an odd prime and $q$ a prime divisor of $a_p$. Then
$q\equiv 1\pmod{4p}$.
\end{proposition}

\begin{proof}
Since $a_p$ is odd, $q$ is odd. The multiplicative order of $2$ modulo $q$
divides $p$ and is not $1$, hence equals $p$. By Fermat's little theorem it
also divides $q-1$, so $p\mid q-1$; as $q-1$ is even and $p$ is odd,
$2p\mid q-1$.
\end{proof}

\begin{remark}
By Zsigmondy's theorem~\cite{Zsi}, for every $n\ge 2$ the number $a_n$ has a
prime divisor that divides none of $a_1,\dots,a_{n-1}$. In view of
Theorem~\ref{thm:main}, such a \emph{primitive} prime divisor of $a_n$
divides no $a_k$ with $k$ not a multiple of $n$.
\end{remark}

\begin{conjecture}\label{conj:mersenne}
There are infinitely many $n$ for which $a_n$ is prime.
\end{conjecture}

\begin{thebibliography}{9}
\bibitem{Luc} \'E.~Lucas, \emph{Th\'eorie des fonctions num\'eriques simplement
p\'eriodiques}, Amer. J. Math. \textbf{1} (1878), 184--240, 289--321.
\bibitem{Zsi} K.~Zsigmondy, \emph{Zur Theorie der Potenzreste}, Monatsh. Math.
Phys. \textbf{3} (1892), 265--284.
\end{thebibliography}

\end{document}
